\documentclass[12pt,a4paper]{article}  % Use this line if this document will be released
%\documentclass[12pt,a4paper,draft]{article}  % Use this line if this document is a draft
%\usepackage{ifdraft}
%% Bibliography
\usepackage{etoolbox}
\newcommand{\bibfile}{\jobname.bib}  % Name of the BibTeX file.
% ref.bib should be a symbolic link to the universal BibTeX file, which should be a local copy of
% https://github.com/equipez/bibliographie/blob/main/ref.bib
% Run `getbib` in the current directory under the draft mode to get the BibTeX file containing only
% the cited references. The name will be xyz.bib if this TeX file is xyz.tex.
\newcommand{\universalbib}{ref.bib}
%\ifdraft{\IfFileExists{\universalbib}{\renewcommand{\bibfile}{\universalbib}}{}}{}
% The counter `cite' is used to count the number of citations.
\newcounter{cite}
\pretocmd{\cite}{\stepcounter{cite}}{}{}


%% Add line numbers in draft mode
\RequirePackage[mathlines]{lineno}
%\ifdraft{\linenumbers}{}
\renewcommand{\linenumberfont}{\normalfont\scriptsize\sffamily\color{gray}}
\setlength{\linenumbersep}{\marginparsep}


%% Geometry
%\voffset=-1.5cm \hoffset=-1.4cm \textwidth=16cm \textheight=22.0cm  % Luis' setting
\usepackage[a4paper, textwidth=16.0cm, textheight=22.0cm]{geometry}
\renewcommand{\baselinestretch}{1.2}


%% Basic packages
\usepackage{amsmath,amsthm,amssymb,amsfonts}
\usepackage{mathtools}  % Provides \coloneqq
\usepackage{empheq}
\usepackage{xcolor}
\usepackage[bbgreekl]{mathbbol}
\DeclareSymbolFontAlphabet{\mathbbm}{bbold}
\DeclareSymbolFontAlphabet{\mathbb}{AMSb}
\usepackage{bbm}
\usepackage{upgreek}
\usepackage{accents}
\usepackage{xspace}
\usepackage{rotating}
\usepackage{multirow,booktabs}
\usepackage[en-US]{datetime2}


%% Format of the table of content
\usepackage[normalem]{ulem}
\usepackage[toc,page]{appendix}
\renewcommand{\appendixpagename}{\Large{Appendix}}
\renewcommand{\appendixname}{Appendix}
\renewcommand{\appendixtocname}{Appendix}
\usepackage{sectsty}
\setcounter{tocdepth}{2}


%% Section title style
\usepackage{sectsty}
\sectionfont{\large}
\subsectionfont{\large}


%% Some colors
\definecolor{darkblue}{rgb}{0,0.1,0.5}
\definecolor{darkgreen}{rgb}{0,0.5,0.1}
\definecolor{darkyellow}{rgb}{0.65,0.65,0.01}


%% Todo notes
%\ifdraft{
%	\setlength{\marginparwidth}{2.42cm}
%	\usepackage[tickmarkheight=3pt,textsize=small,backgroundcolor=blue!16,linecolor=purp%le,bordercolor=purple]{todonotes}
%}{
%	\newcommand{\todo}[1]{}
%	\newcommand{\listoftodos}{}
%}


%% Graph, tikz and pgf
%\usepackage{subfigure}
\setlength{\unitlength}{1mm}
% The \unitlength command is a Length command. It defines the units used in the Picture Environment.
\usepackage{graphicx}
%\usepackage{tikz,tikzscale,pgf,pgfarrows,pgfnodes,filecontents,tikz-cd}
\usepackage{tikz,tikzscale,pgf}
\usetikzlibrary{arrows,arrows.meta,patterns,positioning,decorations.markings,shapes}
\usepackage{pgfplots}
\usepackage{pgfplotstable}
\usepackage[justification=centering]{caption}
\usepgfplotslibrary{fillbetween}
\pgfplotsset{compat=1.11}


%% Turn off some unharmful warnings in draft mode
%% N.B.: DO NOT use `silence` together with `hyperref`. They will cause an infinite loop.
%%\ifdraft{
	%\usepackage{silence}
	%\WarningFilter{xcolor}{Incompatible color definition on}
	%\WarningFilter{hyperref}{Draft mode on}
	%\WarningFilter{refcheck}{Unused label}
	%\WarningFilter{microtype}{`draft' option active}
	%\WarningFilter{latex}{Writing or overwriting file} % Mute the warning about 'writing/overwriting file'
	%\WarningFilter{latex}{Writing file} % Mute the warning about 'writing/overwriting file'
	%\WarningFilter{latex}{Tab has} % Mute the warning about 'Tab has been converted to Blank Space'
%	\WarningFilter{latex}{Marginpar on page} % Mute the warning about 'Marginpar on page xx moved'
	%\WarningFilter{latex}{author given} % Mute the warning about 'No \author given'
%}{}


%% Hyperref, url, and email
%% N.B.: DO NOT use `silence` together with `hyperref`. They will cause an infinite loop.
%\ifdraft{\usepackage{refcheck}\newcommand{\url}{\texttt}}{
	\usepackage{hyperref}
	\hypersetup{colorlinks, linkcolor=darkblue, anchorcolor=darkblue, citecolor=darkblue, urlcolor=darkblue}
	\usepackage{url}
%} % Check unused labels
\newcommand{\email}{\texttt}


%% Enumerate and itemize
\usepackage{eqlist}
\usepackage{enumitem}
\setlist[itemize]{leftmargin=*}
\setlist[enumerate]{leftmargin=*,label=\normalfont{(\alph*)}}


%% Algorithm environment
\usepackage[section]{algorithm}
\usepackage{algpseudocode}
\newcommand{\INPUT}{\textbf{Input}}
\newcommand{\FOR}{\textbf{For}~}
\algrenewcommand\algorithmicrequire{\textbf{Input:}}
\algrenewcommand\algorithmicensure{\textbf{Output:}}
\algrenewcommand\alglinenumber[1]{\normalsize #1.}
\newcommand*\Let[2]{\State #1 $=$ #2}


%% Theorem-like environments
\newtheorem{theorem}{Theorem}[section]
\newtheorem{conjecture}{Conjecture}[section]
\newtheorem{corollary}{Corollary}[section]
\newtheorem{exercise}{Exercise}[section]
\newtheorem{lemma}{Lemma}[section]
\newtheorem{problem}{Problem}[section]
\newtheorem{proposition}{Proposition}[section]
\newtheorem{assumption}{Assumption}[section]
\newtheorem{example}{Example}[section]
\newtheorem{question}{Question}[section]
% Change theoremstyle to ``definition'', which uses textnormal for the text.
\theoremstyle{definition}
\newtheorem{definition}{Definition}[section]
\newtheorem{remark}{Remark}[section]
% proof
\usepackage{xpatch}
\xpatchcmd{\proof}{\itshape}{\normalfont\proofnamefont}{}{}
\newcommand{\proofnamefont}{\bfseries}

%% Equation numbering
\numberwithin{equation}{section}


%% Fine tuning
\usepackage{microtype}
\usepackage[nobottomtitles*]{titlesec} % No section title at the bottom of pages
% Prevent footnote from running to the next page
\interfootnotelinepenalty=10000
% No line break in inline math
\interdisplaylinepenalty=10000
\relpenalty=10000
\binoppenalty=10000
% No widow or orphan lines
\clubpenalty=10000
\widowpenalty=10000
\displaywidowpenalty=10000


% Use @ to put 1 math unit (mu) in math
% See https://nhigham.com/2013/01/07/fine-tuning-spacing-in-latex-equations/
% and also TeXbook p. 155.
\mathcode`@="8000{\catcode`\@=\active\gdef@{\mkern1mu}}


%% Operators, commands
\usepackage{relsize}
\usepackage{nccmath}
%\DeclareMathOperator*{\mcap}{\,\medmath{\bigcap}\,}
%\DeclareMathOperator*{\mcup}{\,\medmath{\bigcup}\,}
\DeclareMathOperator*{\mcap}{\,\mathsmaller{\bigcap}\,}
\DeclareMathOperator*{\mcup}{\,\mathsmaller{\bigcup}\,}
%\renewcommand{\cap}{\mcap}
%\renewcommand{\cup}{\mcup}

\newcommand{\ceil}[1]{ {\lceil{#1}\rceil} }
\newcommand{\floor}[1]{ {\lfloor{#1}\rfloor} }

\DeclareMathOperator{\tr}{tr}
\DeclareMathOperator{\sort}{sort}
\DeclareMathOperator*{\Argmax}{Argmax}
\DeclareMathOperator*{\Argmin}{Argmin}
\DeclareMathOperator*{\Arglocmin}{Arglocmin}
\DeclareMathOperator*{\argmax}{argmax}
\DeclareMathOperator*{\argmin}{argmin}
\DeclareMathOperator*{\diag}{diag}
\DeclareMathOperator*{\Diag}{Diag}
\DeclareMathOperator{\Span}{span}
\DeclareMathOperator{\med}{med}
\DeclareMathOperator{\essinf}{essinf}
\DeclareMathOperator{\cl}{cl}
\DeclareMathOperator{\vol}{vol}
\DeclareMathOperator{\comp}{C}
\DeclareMathOperator{\sign}{sign}
\DeclareMathOperator{\rank}{rank}
\DeclareMathOperator{\range}{range}
\DeclareMathOperator{\card}{card}
\DeclareMathOperator{\diam}{diam}
\DeclareMathOperator{\dist}{dist}
\newcommand{\disth}{{\operatorname{\updelta_{\sss{H}}}}}
\newcommand{\ind}{\mathbbm{1}}
%\newcommand*{\defeq}{\stackrel{\mbox{\normalfont\tiny{\textnormal{def}}}}{=}}
\newcommand\defeq{\mathrel{\overset{\makebox[0pt]{\mbox{\normalfont\tiny\sffamily def}}}{=}}}

\newcommand{\RR}{\mathbb{R}}
\newcommand{\BB}{\mathcal{B}}
\renewcommand{\SS}{\mathbb{S}}
\newcommand{\TT}{\mathcal{T}}
\newcommand{\ZZ}{\mathbb{Z}}
\newcommand{\NN}{\mathbb{N}}
\newcommand{\FF}{\mathcal{F}}
\newcommand{\CC}{\mathbb{C}}
\newcommand{\XX}{\mathcal{X}}
\newcommand{\sset}{\mathcal{S}}
\newcommand{\pen}{h}
\newcommand{\penpar}{\mu}
\newcommand{\res}{\rho}
\newcommand{\col}{r}
\newcommand{\ofd}{\mathcal{F}}
\newcommand{\stf}[1]{\mathbb{S}^{#1}}
\newcommand{\sss}[1]{{\scriptscriptstyle{#1}}}
\newcommand{\sK}{{\scriptscriptstyle{K}}}
\newcommand{\sT}{{\scriptscriptstyle{T}}}
\newcommand{\fro}{{\scriptstyle{\textnormal{F}}}}
\newcommand{\trs}{{\scriptstyle{\mathsf{T}}}}
\newcommand{\hmt}{{\scriptstyle{{\mathsf{H}}}}}
\newcommand{\pin}{{\scriptstyle{{\mathsf{+}}}}}
\newcommand{\inv}{{-1}}
\newcommand{\adj}{*}
\newcommand{\ones}{\mathbf{1}}

\newcommand{\cs}{\text{c}}
\newcommand{\hp}{\circ}
\newcommand{\cc}{\sss{\textnormal{C}}}
\newcommand{\dec}{\sss{\textnormal{D}}}
\newcommand{\cauchy}{\sss{\textnormal{C}}}
\newcommand{\scauchy}{\sss{\textnormal{S}}}
\newcommand{\crit}{\textnormal{crit}}
\newcommand{\rsg}{\hat{\partial}}
\newcommand{\gsg}{\partial}
\newcommand{\dom}{\textnormal{dom}}
\newcommand{\tf}{{\textnormal{f}}}
\newcommand{\tg}{{\textnormal{g}}}
\newcommand{\ts}{{\textnormal{s}}}
\newcommand{\st}{\textnormal{s.t.}}
\newcommand{\etc}{{etc.}\xspace}
\newcommand{\ie}{{i.e.}\xspace}
\newcommand{\eg}{{e.g.}\xspace}
\newcommand{\etal}{{et al.}\xspace}
\newcommand{\iid}{\text{i.i.d.}\xspace}
\newcommand{\as}{\text{a.s.}\xspace}

\newcommand{\me}{\mathrm{e}}
\newcommand{\md}{\mathrm{d}}
\newcommand{\mi}{\mathrm{i}}
\newcommand{\lev}{\mathrm{lev}}
\newcommand{\bA}{\mathbf{A}}
\newcommand{\bx}{\mathbf{u}}
%\newcommand{\bb}{\mathbf{f}}
\newcommand{\bb}{\mathbf{r}}
\newcommand{\nov}{n_{\textnormal{o}}}
\xspaceaddexceptions{]\}}
% tex.stackexchange.com/questions/15252/why-does-xspace-behave-differently-for-parenthesis-vs-braces-brackets
\newcommand{\MATLAB}{\textsc{Matlab}\xspace}
\newcommand{\octave}{\mbox{GNU Octave}\xspace}
\newcommand{\prblm}{\texttt}
\DeclareMathAlphabet{\mathsfit}{T1}{\sfdefault}{\mddefault}{\sldefault}
\SetMathAlphabet{\mathsfit}{bold}{T1}{\sfdefault}{\bfdefault}{\sldefault}
\newcommand{\prbb}{\mathsfit{p}}
\newcommand{\pp}{\mathsf{p}}
\newcommand{\qq}{\mathsf{q}}
\newcommand{\ttt}{\mathsfit{t}}
\newcommand{\tol}{\varepsilon}
\newcommand{\bt}{\mathbf{t}}
\newcommand{\br}{\mathbf{r}}
\newcommand{\dd}{\mathbf{d}}
\newcommand{\ii}{\mathbf{i}}
\newcommand{\jj}{\mathbf{j}}
\newcommand{\xx}{\mathbf{x}}
\renewcommand{\pp}{\mathbf{p}}
\renewcommand{\ggg}{\mathbf{g}}
\newcommand{\GG}{\mathbf{G}}
\DeclareMathOperator{\expc}{\mathbb{E}}
\renewcommand{\Pr}{\mathbb{P}}
\newcommand{\lb}{\underline}
\newcommand{\ub}{\overline}

% mathlcal font
\DeclareFontFamily{U}{dutchcal}{\skewchar\font=45 }
\DeclareFontShape{U}{dutchcal}{m}{n}{<-> s*[1.0] dutchcal-r}{}
\DeclareFontShape{U}{dutchcal}{b}{n}{<-> s*[1.0] dutchcal-b}{}
\DeclareMathAlphabet{\mathlcal}{U}{dutchcal}{m}{n}
\SetMathAlphabet{\mathlcal}{bold}{U}{dutchcal}{b}{n}

% mathscr font (supporting lowercase letters)
%\usepackage[scr=dutchcal]{mathalfa}
%\usepackage[scr=esstix]{mathalfa}
%\usepackage[scr=boondox]{mathalfa}
%\usepackage[scr=boondoxo]{mathalfa}
\usepackage[scr=boondoxupr]{mathalfa}
%\newcommand{\model}{\mathscr{h}}
\newcommand{\model}{\tilde{f}}
\newcommand{\rmod}{F}

\newcommand{\Set}[1]{\mathcal{#1}}
\DeclareMathAlphabet{\mathpzc}{OT1}{pzc}{m}{it} % The mathpzc font
\newcommand{\slv}{\mathpzc}
% mathpzc looks great, but it stops working on 19 Feb 2020 for no reason.
%\newcommand{\slv}{\mathscr}
\newcommand{\software}{\texttt}
\DeclareMathOperator{\eff}{\mathsf{e}\;\!}
\DeclareMathOperator{\Eff}{\mathsf{E}\;\!}
\newcommand{\out}{{\text{out}}}


%% Commands for revision
\newcommand{\red}[1]{\textcolor{red}{#1}}
\newcommand{\blue}[1]{\textcolor{blue}{#1}}
\newcommand{\green}[1]{\textcolor{darkgreen}{#1}}
\newcommand{\TYPO}[1]{{\color{orange}{#1}}}
\newcommand{\MISTAKE}[1]{{\color{violet}{#1}}}
\newcommand{\REPHRASE}[1]{{\color{darkgreen}{#1}}}
\newcommand{\REVISE}[1]{{\color{blue}{#1}}}
\newcommand{\REVISEred}[1]{{\color{red}{#1}}}
\newcommand{\COMMENT}{\todo}  % Needs the todonotes package
%\newcommand{\COMMENT}[1]{\textcolor{brown}{{\small{(comment: #1)}}}}  % This puts comments inline

% Use the following if revision is finished
%\newcommand{\TYPO}{}
%\newcommand{\MISTAKE}{}
%\newcommand{\REPHRASE}{}
%\newcommand{\REVISE}{}
%\newcommand{\REVISEred}{}
%\newcommand{\COMMENT}[1]{}  % Input ignored.

\renewcommand{\algorithmicrequire}{\textbf{Input:}} % Redefine Input
\renewcommand{\algorithmicensure}{\textbf{Output:}}
%%%%%%%%%%%%%%%%%%%%%%%%%%%%%%%%%%%%%%%%%%%%%%%%%%%%%%%%%%%%%%%%%%%%%%%%%%%%%%%%%%%%%%%%%%%%%%%%%%%%
\title{ Origin, Dilemma and Development of the DFP Algorithm}

\date{\DTMnow}

\author{Chen Huicong}	
	%\and
	%Author2
	%\thanks{Information2}

\begin{document}

\maketitle

%\begin{abstract}
%\end{abstract}

%\textbf{Keywords}: Keyword1, Keyword2
%%%%%%%%%%%%%%%%%%%%%%%%%%%%%%%%%%%%%%%%%%%%%%%%%%%%%%%%%%%%%%%%%%%%%%%%%%%%%%%%%%%%%%%%%%%%%%%%%%%%



%%%%%%%%%%%%%%%%%%%%%%%%%%%%%%%%%%%%%%%%%%%%%%%%%%%%%%%%%%%%%%%%%%%%%%%%%%%%%%%%%%%%%%%%%%%%%%%%%%%%
%% References
% Include references only if there are citations.

\section{Introduction to the DFP algorithm}
The DFP algorithm is the first quasi-Newton method for unconstrained optimization
\begin{equation}
	\min_{x \in \mathbb{R}^n} f(x)\text{.}
\end{equation}
Davidon first proposed this algorithm in 1959\cite{doi:10.1137/0801001}. Fletcher and Powell then systematically organized and revised it in 1963\cite{10.1093/comjnl/6.2.163}.



\subsection*{\normalsize Algorithm 1.1: The Davidon-Fletcher-Powell algorithm}\label{1.1}

\begin{enumerate}[leftmargin=4em]
	\item[Step 0] Given initial point  $x_0 \in 
	\RR^n$; matrix $B_0\in\RR^{n\times n}$ (commonly $B_0= I$)~positive definite;  convergence tolerance $\epsilon > 0$. Set $k=0$. 
	
	\item[Step 1] Compute $g_k=\nabla f(x_k)$;\\ 
	if $\|g_k\|~\leq~\epsilon$, then stop;\\
	set $p_k=-B_k^{-1}g_k$. 
	
	\item[Step 2]Carry out a line search along $p_k$, getting $~\alpha_k > 0$\\
	satisfying weak Wolfe conditions, \begin{subequations}
		\label{eq:group}        
		\begin{align}
			f({x_k} + \alpha_k {p_k}) &\leq f({x_k}) + c_1 \alpha_k \nabla f({x_k})^\top {p_k} \label{eq:wolfe1} \\
			\nabla f({x_k} + \alpha_k {p_k})^\top {p_k} &\geq c_2 \nabla f({x_k})^\top {p_k} \label{eq:wolfe2}
		\end{align}
	\end{subequations}
	where $0 < c_1 < c_2 < 1$~(commonly $c_1=10^{-4}$, $c_2=0.9$).\\
	Set $~x_{k+1}=x_k+\alpha_k p_k$. 
	
	\item[Step 3](DFP update formula) Set
	\begin{equation}
	B_{k+1} = B_k - \frac{B_k s_k y_k^\top + y_k s_k^\top B_k}{y_k^\top s_k}
	+ \left(1 + \frac{s_k^\top B_k s_k}{y_k^\top s_k}\right) \frac{y_k y_k^\top}{y_k^\top s_k}.\label{1.3}
    \end{equation}
	where\begin{equation}
		s_k=\alpha_k p_k
	\end{equation}
	\begin{equation}
	y_k=g_{k+1}-g_k
    \end{equation}
	\item[Step 4] $~k=k+1$; go to Step 1. \\
	
\end{enumerate}
The BFGS method  proposed by Broyden, Fletcher, Goldfarb, and Shanno \cite{10.1093/imamat/6.1.76}, is the same as formula \eqref{1.3}, which is replaced by
\begin{equation}
	B_{k+1} = B_k - \frac{B_k s_k s_k^\top B_k}{s_k^\top B_k s_k} + \frac{y_k y_k^\top}{y_k^\top s_k}.
\end{equation}
\section{Superlinear Convergence of the DFP algorithm}
Powell\cite{powell1971convergence} showed that for uniformly convex functions, the DFP algorithm with exact line search is globally convergent.
 Dennis\cite{636c0025-00c1-3c96-8fa3-324390e0f75f} proposed a necessary and sufficient condition for the superlinear convergence of quasi-Newton methods. 
\begin{theorem}
	Let \(F:\mathbb{R}^n \to \mathbb{R}^n\) be differentiable in an open convex set \(D \subseteq \mathbb{R}^n\), and assume that for some \(x^* \in D\), \(F'\) is continuous at \(x^*\) and \(F'(x^*)\) is nonsingular. Let \(\{B_k\} \subset \mathbb{R}^{n \times n}\) be a sequence of nonsingular matrices, and suppose that for some \(x_0 \in D\) the sequence defined by
	\begin{equation}
	x_{k+1} = x_k - B_k^{-1} F(x_k), \qquad k = 0,1,2,\dots
	\end{equation}
	remains in \(D\) and converges to \(x^*\). Then \(\{x_k\}\) converges \(Q\)-superlinearly to \(x^*\) and \(F(x^*) = 0\) if and only if
	\begin{equation}
	\lim_{k\to\infty} \frac{\big\| [B_k - F'(x^*)] (x_{k+1} - x_k) \big\|}{\| x_{k+1} - x_k \|} = 0.
	\end{equation}
\end{theorem}
Based on this theorem, we can obtain some conditions on the superlinear convergence of the DFP algorithm or other quasi-Newton methods.\cite{d4c4e1e1-b552-312d-a127-ddfcff7c0c72}

\section{The Global Convergence of DFP algorithm}
Powell\cite{powell1971convergence} proved that if $f(x)$ is a convex function and if $\alpha_k~$in step 2 satisfies
\begin{equation}
	~\alpha_k~=~\argmin_{\alpha>0}f(x_k+\alpha p_k).
\end{equation}
That is, the line searches in step 2 are exact in all iterations. Then the DFP method is globally convergent. For the inexact line search, Powell\cite{zbMATH03529352} proved that the BFGS method is globally convergent. Then Byrd, Nocedal and Yuan\cite{badf1033-31a5-3157-8131-6cdeb5b73b0c} proved the global convergence of all methods in the convex Broyden family except the DFP method. But for DFP method with inexact line searches, we still do not have an exact result.

Yuan\cite{Yuan1995convergence} introduced a new method for analyzing the square trace of matrices and presented some sufficient conditions for the global convergence of DFP algorithm when it satisfies

\begin{assumption}\label{assum}
	\item[(i)] the objective function $f(x)$ is twice continuously differentiable and uniformly convex, i.e.~, there exists a constant \(c_3 > 0\) such that
	\begin{equation}
		p^{\top}\nabla^2 f(x)p \ge c_3 \| p \|_2^2,\quad \forall p,x\in\mathbb{R}^n.
	\end{equation}
	\item[(ii)] In all iterations, the line search satisfies the weak Wolfe conditions. 
\end{assumption}  
Yuan proved the following four theorems as conditions for DFP algorithm convergence under the assumption \ref{assum}.
\begin{theorem}
	 Under the assumption \ref{assum}, if the step sizes \(\alpha_k\) are bounded away from zero (i.e., there exists \(\varepsilon > 0\) such that \(\alpha_k \ge \varepsilon\) for all \(k\)) and the line search satisfies the strong Wolfe conditions. That is, keep \eqref{eq:wolfe1} and change \eqref{eq:wolfe2} to
	\begin{equation}
	|p_k^{\top} \nabla f(x_k + \alpha_k p_k)| \le -c_2 p_k^{\top} g_k.
	\end{equation}
	Then the sequence \(\{x_k\}\) generated by the DFP method converges to the unique minimizer of \(f\).
\end{theorem}


\begin{theorem}
	 Under the assumption \ref{assum}, if the sequence \(\{x_k\}\) generated by the DFP algorithm satisfies \begin{equation}
	 	\|g_{k+1}\|_2 \le \|g_k\|_2
	 \end{equation} for all sufficiently large \(k\), then \(\{x_k\}\) converges to the unique minimizer of \(f\).
\end{theorem}


\begin{theorem}
	Under the assumption \ref{assum}, if the displacement vectors \(s_k = x_{k+1} - x_k\) satisfy
	\begin{equation}
	\sum_{k=1}^{\infty} \|s_k\| < \infty,
	\end{equation}
	then the DFP algorithm converges to the solution.
\end{theorem}

\begin{theorem}
	Under the assumption \ref{assum}, if \begin{equation}
		y_k^T B_k g_k \leq 0, 
	\end{equation} for all \(k\), then \(\{x_k\}\) either terminates at the unique minimum \(~x^*\) of \(~f(x)\) or converges to \(~x^*\).
\end{theorem}
At the end of the paper, Yuan also proposed two inequalities as open problems, either of which could equally serve as a condition for the convergence of the DFP algorithm.

However, since 1995, there has been no major progress on the DFP algorithm. Although many sufficient conditions\cite{YYSU199802003,GXKZ200201015,DQSY200101020} have been proposed over the past 30 years, the global convergence of the DFP algorithm has still not been fully resolved, and it remains a research direction in the field of quasi-Newton methods. 

\section{Improvement of the DFP Algorithm}
Due to the better stability and global convergence of other algorithms in the Broyden family (such as the BFGS algorithm), people prefer to use them over the DFP algorithm. Additionally, attempts have been made to improve the DFP algorithm to achieve better convergence, which represents the current development direction of the DFP algorithm.

Gonglin Yuan, Peifeng Zhou and Hongtruong Pham\cite{yuan2026projection} propose Projection DFP quasi-Newton Algorithm (P-DFP-NA), which is an effective improvement of the DFP algorithm. The following $~p_k=-B_k^{-1}g_k$ and $\alpha_k$ are the same with Algorithm \hyperref[1.1]{1.1}. Then we define a set by \begin{equation}
	~PX_k = \{x_k \mid g_k^T p_k \leq -\tau \|p_k\|\},
\end{equation} where $~\tau ~> ~0$ is a constant. The iteration points in $PX_k$ are in-points, and other points that do not belong to the set $PX_k$ are out-points. We discuss the following two cases.
\textbf{Case i} The iteration point $x_k$ is in-point. Then $x_{k+1}=x_k+\alpha_k p_k$, which is the same as the DFP method.\\
\textbf{\normalsize Case ii} The iteration point $x_k$ is out-point. We define $h_k=x_k+\alpha_kp_k$, and the parabolic surface \begin{equation}
	\{x \mid g(h_k)^T p_k - \sigma_1 g(x)^T p_k + \sigma_2 \tau \|p_k\| = 0\}
\end{equation}
where \( \sigma_1 + \sigma_2 < 1 \) with \( \sigma_1 \in (0, c_2) \) and \( \sigma_2 \in (c_2 - \sigma_1, 1) \) are constants, and \(c_2\) is the parameter in the weak Wolfe conditions \eqref{eq:wolfe2}. Then \begin{equation}
	x_{k+1} := x_k + \frac{PP_k}{\|p_k\|} \alpha_k p_k.
\end{equation}
where \( PP_k = g(h_k)^T p_k - \sigma_1 g(x_k)^T p_k + \sigma_2 \tau \|p_k\| \). We can prove that \begin{theorem}
	In Projection DFP algorithm, $s_k=x_{k+1}-x_k$ satisfies
	\begin{equation}
		\sum_{k=1}^{\infty} \|s_k\| < \infty.
	\end{equation}
\end{theorem}
Furthermore, we can prove that this algorithm globally converges.
\ifnum\value{cite}>0
    \small
    \bibliographystyle{plain}
    \bibliography{ref.bib}
\fi
%% The end

\end{document}